\chapter{What rosetta is}
\label{ch:what}

\section{The problem, stated once}

Exposing a C++ library to another language is normally paid for once per
destination. A pybind11 module for Python. An N-API layer for Node. A sol2
usertype for Lua. A P/Invoke shim for C\#. A moc-annotated hierarchy for a Qt
panel. A JSON route table for a REST service.

Look at what those files actually contain. Here is one method, five times:

\begin{lstlisting}[style=cpp]
double Grid::at(int i, int j) const;
\end{lstlisting}

\begin{lstlisting}[style=cpp]
py::class_<Grid>(m, "Grid").def("at", &Grid::at);                 // pybind11
ut["at"] = &Grid::at;                                             // sol2
exports.Set("at", Napi::Function::New(env, wrap_at));             // N-API
emscripten::class_<Grid>("Grid").function("at", &Grid::at);        // embind
[DllImport("grid")] static extern double Grid_at(IntPtr h, int i, int j);
\end{lstlisting}

Every one of those restates, in a different notation, something the C++
compiler already knows: that \code{Grid} has a member called \code{at}, that it
takes two \code{int}s, that it returns a \code{double}, that it does not
mutate. The information is not missing. It is simply not reachable from
anywhere except the compiler.

And the cost is not the initial write --- it is the \emph{derivative}. Add a
parameter to \code{at} and five files are now wrong, in five languages, and
four of them will not tell you until run time.

\section{The idea}

C++26 static reflection (P2996) makes that information reachable from inside
C++ itself. A program can ask, at its own compile time, what members a type
has, what their types are, whether they are public, whether they are virtual.

Rosetta does exactly one thing with that: it reads your headers \emph{once},
through reflection, into a target-neutral intermediate representation, and then
\emph{projects} that representation into as many bindings as you asked for.

\begin{center}
\begin{tikzpicture}[
  node distance=6mm,
  box/.style={draw,rounded corners=2pt,align=center,font=\small,inner sep=5pt,minimum height=8mm},
  user/.style={box,fill=blue!6,draw=rosettablue},
  tool/.style={box,fill=orange!8,draw=rosettaorange},
  outbox/.style={box,fill=black!4,draw=black!35,dashed},
  arr/.style={-{Stealth[length=2mm]},draw=black!55}
]
\node[user] (h) {your headers\\\scriptsize unmodified};
\node[user,below=of h] (m) {\texttt{manifest.json}\\\scriptsize what to bind};
\node[tool,right=14mm of $(h)!0.5!(m)$] (r) {reflection\\walk};
\node[outbox,right=12mm of r] (ir) {IR\\\scriptsize target-neutral};
\node[outbox,above right=2mm and 12mm of ir] (a) {\texttt{bindings/python}};
\node[outbox,right=12mm of ir] (b) {\texttt{bindings/node}};
\node[outbox,below right=2mm and 12mm of ir] (c) {\texttt{bindings/\ldots}};
\draw[arr] (h) -- (r);
\draw[arr] (m) -- (r);
\draw[arr] (r) -- (ir);
\draw[arr] (ir) -- (a);
\draw[arr] (ir) -- (b);
\draw[arr] (ir) -- (c);
\end{tikzpicture}
\end{center}

You declare \emph{what} to bind. You never declare \emph{how}.

\section{What you write}

Two things, and one of them you already have.

The header stays exactly as it is:

\begin{lstlisting}[style=cpp]
// my_lib/person.h
#include <string>

struct Person {
    std::string name;
    int         age = 0;
    std::string id;
    std::string greet(const std::string &salutation) const;
};
\end{lstlisting}

The manifest is the only new file:

\begin{lstlisting}[style=json]
{
  "user_include": "../my_lib",
  "rosetta_include": "/path/to/rosetta/include",
  "targets": [
    { "lang": "python", "name": "person_py" },
    { "lang": "node",   "name": "person_js" }
  ],
  "classes": [
    { "name": "Person", "header": "person.h" }
  ]
}
\end{lstlisting}

Then one command:

\begin{lstlisting}[style=sh]
rosetta_gen --build manifest.json
\end{lstlisting}

and \code{bindings/python} and \code{bindings/node} exist, compiled, with
\code{Person} in both --- its three fields as properties, its method as a
method, its constructors as constructors.

\section{What comes out without being asked for}

The manifest above names one class. Reflection supplies the rest of it:

\begin{itemize}
  \item \textbf{Public fields}, as read/write properties with per-backend
    getters and setters.
  \item \textbf{Public methods}, instance and \code{static} alike.
  \item \textbf{Constructors} --- every overload, not just the default one.
  \item \textbf{Inheritance}. Public base fields and methods are flattened into
    the derived binding; a derived declaration shadows the base one, and a
    virtual diamond collapses to a single member. Virtual methods are flagged,
    so backends that can generate trampolines do.
  \item \textbf{Overload sets}. The whole set reaches the IR; what each target
    does with it depends on whether its runtime can dispatch on argument types
    (Chapter~\ref{ch:coverage}).
  \item \textbf{Enums}, both plain and scoped, with their enumerators.
  \item \textbf{Nested user types and \code{std::vector}}, marshalled across
    the boundary --- a \code{Surface} returning \code{std::vector<Triangle>}
    works without a line about \code{Triangle} beyond listing it.
\end{itemize}

Free functions are the one thing you must name, because a namespace has no
member list to enumerate --- you list them in the manifest, still without
touching the header (Section~\ref{sec:functions}).

\section{What you get out}

One self-contained project per target, under \code{bindings/}:

\begin{lstlisting}[style=plain]
bindings/
  python/    auto_pybind.cpp      CMakeLists.txt  pyproject.toml  make_wheel.py
  node/      auto_napi.cpp        CMakeLists.txt  package.json
  wasm/      auto_emscripten.cpp  CMakeLists.txt
  typescript/ person.d.ts
  coverage.json
\end{lstlisting}

Each is an ordinary CMake (or npm, or emcmake) project. Nothing in it refers to
rosetta at run time; nothing in it needs reflection to compile. That last point
is worth its own section.

\section{Reflection-free by construction}
\label{sec:reflection-free}

Every compiled backend \emph{fully expands} what it found. There are no splices
left in the output, no \code{<experimental/meta>}, no trace that reflection was
ever involved --- just explicit pybind11 / N-API / embind / sol2 / jlcxx calls,
or, for the \lang{dynamic} target, aggregate-initialised tables.

The consequence is the one that makes rosetta usable in practice:

\begin{note}
Reflection runs \emph{once}, on the generation host. The machine that builds
the generated binding needs an off-the-shelf toolchain --- plain Clang, GCC,
MSVC, a stock emsdk, a stock Qt~6. It does not need a C++26 compiler, and it
does not need rosetta.
\end{note}

So the workflow for a project whose CI runs on a normal compiler is: generate
on one machine that has the fork, commit or ship the generated tree, and build
it anywhere.

Two targets are exceptions and are marked as such throughout: \lang{rest} and
\lang{json} emit code that still splices reflections, so those two --- and only
those two --- need the C++26 toolchain at their own build time.

\section{The nineteen targets}

\begin{center}
\small
\begin{tabular}{@{}llL{5.4cm}@{}}
\toprule
\textbf{Kind} & \textbf{\code{lang}} & \textbf{Output} \\
\midrule
Scripting  & \lang{python}     & pybind11 extension module \\
           & \lang{nanobind}   & nanobind extension module \\
           & \lang{node}       & N-API native addon \\
           & \lang{wasm}       & Emscripten / embind module \\
           & \lang{lua}        & sol2 module, \code{require}-able \\
           & \lang{julia}      & CxxWrap.jl / jlcxx shared module \\
\addlinespace
Managed    & \lang{csharp}     & C-ABI shared library + P/Invoke wrappers \\
           & \lang{java}       & C-ABI + handle-backed FFM wrappers \\
\addlinespace
User inter- & \lang{qt}        & Qt Widgets property/method inspector \\
faces      & \lang{qml}        & QML / QtQuick inspector \\
           & \lang{imgui}      & Dear ImGui inspector application \\
\addlinespace
Services   & \lang{rest}       & cpp-httplib JSON server + browser client \\
           & \lang{openapi}    & OpenAPI 3.1 specification \\
           & \lang{json}       & nlohmann (de)serialisation \\
\addlinespace
Documents  & \lang{typescript} & ambient \code{.d.ts} declarations \\
           & \lang{markdown}   & API reference document \\
           & \lang{html}       & styled API reference page \\
           & \lang{paraview}   & ParaView Server Manager plugin XML \\
\addlinespace
Metadata   & \lang{dynamic}    & the IR as runtime data (Chapter~\ref{ch:dynamic}) \\
\bottomrule
\end{tabular}
\end{center}

The full table, with what each one needs to build, is
Appendix~\ref{app:targets}. The last row is not a language binding at all and
is the subject of its own chapter.

\section{Annotations are optional}

Reflection binds the class. \emph{Annotations} enrich it --- a docstring, a
numeric range enforced at the boundary, a read-only property, a UI hint that
turns a field into a slider:

\begin{lstlisting}[style=cpp]
[[= rosetta::doc{"Max solver iterations"}, = rosetta::range{0, 200},
   = rosetta::widget::slider ]]
int maxIter{100};
\end{lstlisting}

They are never required, they never change the C++ behaviour of your type, and
you add them only where they pay off. If you would rather not touch the header
at all --- because it is a third party's, or because you want the generated
binding to build on a stock compiler --- the same metadata goes in a JSON
side-car instead, wired by one manifest field. Chapter~\ref{ch:annotations}
covers both forms.

\section{What rosetta is not}

Being clear about this early saves disappointment later.

\begin{itemize}
  \item \textbf{It is not a C++ parser.} It does not read your headers as text.
    A compiler does, and rosetta asks the compiler. The one place a heuristic
    scan exists --- \code{rosetta\_gen --init} over a source tree --- is
    explicitly a starting point you then edit.
  \item \textbf{It is not a runtime.} Nothing rosetta ships is loaded by your
    finished binding. The generated code depends on pybind11, or on N-API, or
    on sol2 --- not on rosetta.
  \item \textbf{It does not bind everything.} Some C++ has no shape in the
    destination language: a \code{std::function} parameter, a raw pointer
    return, an operator with no name to bind to. Those members are
    \emph{skipped}, never fatal, and every skip is recorded with a reason in
    \code{coverage.json} (Chapter~\ref{ch:coverage}). A binding that silently
    loses a method is the failure mode this design exists to avoid.
  \item \textbf{It is not finished.} It is a prototype tracking proposals that
    are themselves not final. See the status note in the Preface.
\end{itemize}

\section{Where to go next}

Chapter~\ref{ch:pipeline} explains why generating a binding is a
\emph{pipeline} rather than a command --- which is the one piece of rosetta's
shape that is not obvious, and the piece that explains most of the rest.
If you would rather see it work first, skip to Chapter~\ref{ch:first}.
